% !TeX root = ./main.tex

% --------------------------------------------------
% 資訊設定（Information Configs）
% --------------------------------------------------

\ntusetup{
  university*   = {National Taiwan University},
  university    = {國立臺灣大學},
  college       = {生物資源暨農學院},
  college*      = {College of Bioresources and Agriculture},
  institute     = {生物機電工程學系},
  institute*    = {Department of Biomechatronics Engineering}, % 請確認您的科系為Department of xxx或Institute of xxx
  title         = {論文標題},
  title*        = {Title},
  author        = {論文作者},
  author*       = {Name of author},
  ID            = {學號},
  advisor       = {吳筱梅},
  advisor*      = {Hsiao-Mei Wu},
  date          = {2020-05-01},         % 若註解掉，則預設為當天
  oral-date     = {2020-05-01},         % 若註解掉，則預設為當天
  DOI           = {10.5566/NTU2018XXXXX},
  keywords      = {關鍵字1, 關鍵字2, 關鍵字3, 關鍵字4},
  keywords*     = {Keywords1, Keywords2, Keywords3, Keywords4},
}

% --------------------------------------------------
% 加載套件（Include Packages）
% --------------------------------------------------

\usepackage[authoryear,round]{natbib}   % 參考文獻
\usepackage{amsmath, amsthm, amssymb}   % 數學環境
\usepackage{ulem, CJKulem}              % 下劃線、雙下劃線與波浪紋效果
\usepackage{booktabs}                   % 改善表格設置
\usepackage{multirow}                   % 合併儲存格
\usepackage{diagbox}                    % 插入表格反斜線
\usepackage{array}                      % 調整表格高度
\usepackage{longtable}                  % 支援跨頁長表格
\usepackage{paralist}                   % 列表環境


\usepackage{lipsum}                     % 英文亂字
\usepackage{zhlipsum}                   % 中文亂字

% --------------------------------------------------
% 套件設定（Packages Settings）
% --------------------------------------------------

\makeatletter
\newcounter{bibcount}
\setlength{\bibhang}{3em}
\let\old@bibitem\bibitem
\renewcommand{\bibitem}[2][]{%
  \old@bibitem[#1]{#2}%
  \stepcounter{bibcount}%
  \makebox[\bibhang][l]{[\thebibcount]}\ignorespaces
}
\makeatother
